\section{Synthesis and Limitations}

The frozen Qwen3.8 judge retains the 40-example seed because the optimized
candidate fails the per-class recall gate. Its NVFP4 deployment reaches
0.710 accuracy on 407 reused validation units under both low and medium
effort, while changing 10.1\% of labels between settings. Neither prompt
selection nor this deployment check establishes annotation accuracy on the
generated benchmark traces. Gemma and GPT-OSS now share a stratified 100,000-sentence sampling
protocol and judge configuration, with 99,999 and 100,000 valid labels
respectively. Their selected cohorts remain model-specific; Qwen uses a
different sampling protocol and judge deployment.

All three models retain prospective episode information on held-out
responses from the released SAT corpus. Five response-disjoint outer folds
with inner-validation layer selection give macro-AUROC 0.852 ([0.838,0.868])
for Qwen, 0.780 ([0.761,0.797]) for GPT-OSS, and 0.788 ([0.772,0.805]) for Gemma.
Each exceeds its two-feature prefix-position baseline. These 95\% intervals
resample responses conditional on the saved folds and fits. The study
strengthens same-corpus generalization evidence but leaves external transfer,
richer shortcut controls, and online detector validity unresolved.
The original sentence-split scores and routing-layer choices retain their
overlap and test-selection limitations. Qwen's external 581-unit inspection has 50.3\%
primary and 60.8\% compatible agreement under a provisional hybrid review;
it is neither independent human validation nor evidence for the other encoders.

The three correlation corpora each cover 4,647 attempts, 1,547 problems,
and 100 complete avg@32 groups. Associations depend on model, benchmark,
and aggregation level. For Qwen, whole-continuation confidence correlates
with avg@32 at $\rho=0.632$ on AIME24 and $0.705$ on AIME25,
whereas AMC23 gives $-0.200$ with an interval containing zero.
These estimates do not control length, termination, or difficulty.
Problem-respecting bootstrap intervals address repeated attempts for the
reported primary features; naive attempt-level significance screens remain
exploratory. Class coverage now spans all six benchmarks for Gemma and GPT-OSS.
Sparse class/outcome cells and unmatched problem cohorts still limit
comparisons; equal sentence budgets do not balance benchmark or class shares.
No common causal interpretation of router confidence follows from these data.

The stratified class analysis also shows benchmark-dependent signs. For
Gemma's Implement-labelled tokens, confidence correlates with correctness
at $r=0.498$ on AIME24 (30 attempts) and $-0.367$ on AMC23 (40 attempts).
For GPT-OSS, the corresponding coefficients are $0.164$ and $-0.098$.
These are descriptive estimates on different model-specific cohorts,
not paired tests of model differences. Gemma's AMC23 Plan, Explore, and
Monitor cohorts contain only correct attempts, so their correctness
correlations are undefined despite nonzero label coverage.

After consistent offline rescoring, no complete intervention establishes
a beneficial accuracy direction using the reported 95\% interval criterion.
Qwen fixed positive strength $+2$ loses $2.80$ attempt-weighted points
($[-5.16,-0.72]$); GPT-OSS fixed positive strength $+2$ loses $3.30$
($[-6.16,-0.93]$). GPT-OSS paper-style suppression loses $2.22$ points
($[-4.16,-0.43]$), and positive/negative paper-style promotion loses
$61.65$ and $59.57$ points with intervals well below zero.
All stochastic margin intervals contain zero: Qwen concurrency 25 gives
$+0.93$ points ($[0.00,+1.86]$), while GPT-OSS gives $-0.86$
($[-2.87,+1.15]$). Equal-problem weighting removes the below-zero
intervals for fixed strength $+2$ and paper-style suppression, but
paper-style promotion remains strongly negative. The earlier positive
OSS secondary intervals and greedy Qwen margin interval do not survive
rescoring. These exploratory intervals remain conditional on saved
attempts and frozen policies and are not adjusted for multiple comparisons.

The Qwen paper-style positive intervention was interrupted after 494 of
1,395 attempts. On those same completed IDs, baseline accuracy is 70.45\%
and guided accuracy is 0.81\%; all 494 guided outputs hit the 32,768-token
limit. This subset covers all AIME24 attempts and part of AIME25 only,
so it does not estimate the full evaluation effect
(Section~\ref{sec:partial-qwen-paper}). Each effect uses
its own matched baseline, generated or reused after compatibility and
output-hash checks. Calibration excludes unscored traces and selects targets
from observational associations; Qwen also transfers across quantizations.
Margin manipulation changes expert membership in finite precision. Strong
paper steering retains our selection estimator and is not a full replication
of the paper. The evaluated controllers do not use online episode
predictions: annotations and probes support measurement and target exploration.
The qualitative sample in Section~\ref{sec:paper-failure-inspection}
contains direct repetition examples in capped and normally ended outputs,
an arithmetic error, and two equivalent final answers rejected by the saved
scorer. These descriptions are conditional on a small stratified sample;
they establish neither population failure-mode frequencies nor a causal
mechanism. The documented formatting discrepancies motivated the completed offline
rescoring pass; the 22-case inspection retains its original sampling frame.
These findings test the actionability of routing associations, without
establishing a portable policy or isolating identity from mixture-weight
changes. Intervals are conditional on saved runs, omit independent seed
variation, and are not adjusted for multiple comparisons. Objective and
strength selection require an independent evaluation design.


The full-diagnostics GPT-OSS rerun additionally tests a fixed-cache proxy
(Section~\ref{sec:oss-cache-diagnostics}). Target-expert coverage increases
from $5.08\%$ to $35.17\%$, with a rescored accuracy change of $+0.14$
points ($[-2.37,+2.37]$). Yet a calibration-frequency cache already covers
$36.87\%$ before bias, and guided outputs are $4.88\%$ longer on average.
Greater target reuse therefore does not by itself establish fewer expert
loads or faster inference. A modest accuracy loss may be useful if offset
by sufficient speed, but the accuracy--latency trade-off remains unmeasured.



The completed Gemma fixed-bias $+1$ evaluation gives $87.67\%$ baseline
and $87.60\%$ guided accuracy, a $-0.07$-point change
($[-1.15,+1.00]$). Across both diagnostic sessions, target coverage rises
from $0.64\%$ to $5.03\%$, still far below the $41.81\%$ pre-bias
frequency-cache coverage. The interrupted-session counters include
execution work beyond the final scored outputs.

The matched 100-problem ordered-cache experiments
(Section~\ref{sec:multimodel-cache-100}) reveal model-dependent effects.
Rescored accuracy changes are $-1$ point for OSS and Qwen and zero for
Gemma, with all intervals including zero. Qwen reduces loads per token
by $5.20$--$6.82\%$, with intervals below zero, and has promising
$6.20$--$7.81\%$ reductions in observed loads per answer whose intervals
still include zero. Gemma target pinning increases loads per answer by
$10.32\%$ at 16 slots ($[+2.69,+18.68]\%$), while guided LRU avoids most
of that cost. OSS's longer outputs offset its 16-slot per-token savings.
These results distinguish routing locality, pinning policy and total
answer cost; none measures actual SSD traffic or offloaded-inference speed.
